% ===================================================================
\section{Core Verification Results}
\label{ch:core}
% ===================================================================

\subsection{Module Architecture}

The 21 core modules are organized in five functional layers:

\begin{center}
\begin{tabular}{lll}
\toprule
Layer & Modules & Role \\
\midrule
Kernel (L1) & KernelDefinitions & Inductive types, primitives, objects \\
Float (L2) & FloatTheory, FloatMonotonicity & IEEE 754 arithmetic foundations \\
Structure (L3) & CategoryStructure, PipelineComposition, & \\
 & GrothConstruction, Termination & Core categorical and operational structure \\
Semantics (L4) & SemanticCorrectness, TacitDependence, & \\
 & ExplicitArtifact, ReductionSemantics, & \\
 & ConfluenceProof, RunsTrajectories & Proof semantics and equivalence \\
Governance (L5) & CollectiveLayer, GovernanceProofs, & \\
 & AmbiguityProtocol, CredenceFramework, & \\
 & DegradationModel, StoppingTime, & \\
 & SurvivalAnalysis, GovernanceSafeExtension & Governance and probabilistic layer \\
\bottomrule
\end{tabular}
\end{center}

% -------------------------------------------------------------------
\subsection{L1: Kernel Definitions}
% -------------------------------------------------------------------

\leanref{Frfp/Core/KernelDefinitions.lean}

The kernel establishes the three FRFP objects and the closed 6-element primitive
set.

\begin{lstlisting}
inductive Object where
  | Explicit | Tacit | Boundary
  deriving DecidableEq, Repr

inductive Primitive where
  | RunInfo | ExecutionContext | ExecutionDecision
  | RunBoundary | TacitElement | HumanFeedbackData
  deriving DecidableEq, Repr
\end{lstlisting}

\textbf{Key theorems (8):}
\begin{itemize}
  \item \texttt{ets\_unique}: Explicit-tacit separation is the unique non-degenerate
        partition of \texttt{Object} satisfying separation axioms.
  \item \texttt{rb\_exactly\_one}: Run-boundary has exactly one well-typed
        authorization witness per protocol step.
  \item \texttt{ai\_executable}: All and only explicit-region objects satisfy the
        AI-executability predicate.
\end{itemize}

% -------------------------------------------------------------------
\subsection{L2: Float Theory and Monotonicity}
% -------------------------------------------------------------------

\leanref{Frfp/Core/FloatTheory.lean}, \leanref{Frfp/Core/FloatMonotonicity.lean}

These modules provide the 59 IEEE 754 axioms and prove 31 theorems about
floating-point credence arithmetic, including:

\begin{itemize}
  \item \texttt{credence\_clamp\_range}: Credence values are bounded in $[0.0,\,1.0]$
        after clamping.
  \item \texttt{float\_foldl\_monotone}: Folding a monotone function preserves
        order.
  \item \texttt{credence\_product\_le\_min}: Products of credences are bounded by
        the minimum operand.
\end{itemize}

\begin{lstlisting}
theorem credence_product_le_min (x y : Float) (hx : 0 ≤ x) (hy : 0 ≤ y)
    (hx1 : x ≤ 1) (hy1 : y ≤ 1) : x * y ≤ min x y := by
  simp [min_def]
  split
  · exact float_mul_le_of_le_one_right hx hy1
  · exact float_mul_le_of_le_one_left hy hx1
\end{lstlisting}

% -------------------------------------------------------------------
\subsection{L3: Structure Layer}
% -------------------------------------------------------------------

\leanref{Frfp/Core/CategoryStructure.lean}

The category structure module proves that \(\Ecat\) and \(\Tcat\) satisfy the
Mac~Lane category axioms (identity, associativity). The Grothendieck construction
\(\int F\) is built as a dependent product over the base category.

\leanref{Frfp/Core/ConfluenceProof.lean}

The confluence module uses a three-step argument:
\begin{enumerate}
  \item Prove well-founded termination of the reduction relation via a step-count
        measure.
  \item Prove local confluence (diamond property) by case analysis over the 18
        reduction rules.
  \item Invoke Newman's Lemma: termination + local confluence implies global
        confluence.
\end{enumerate}

\begin{lstlisting}
-- Newman's Lemma application
theorem global_confluence : Confluent frfp_reduce := by
  apply NewmansLemma.confluent_of_wf_locally_confluent
  · exact termination_wf
  · exact local_confluence
\end{lstlisting}

% -------------------------------------------------------------------
\subsection{L4: Semantics Layer}
% -------------------------------------------------------------------

\leanref{Frfp/Core/SemanticCorrectness.lean}

The semantics layer proves:
\begin{itemize}
  \item \texttt{semantic\_correctness}: The denotational semantic function
        \(\mathrm{Sem}\) is well-defined and type-correct.
  \item \texttt{observation\_invariance}: Observable output is invariant under
        navigation-equivalent terms.
  \item \texttt{pipeline\_immutability}: Pipeline semantics are stable under
        authorized composition.
\end{itemize}

\leanref{Frfp/Core/TacitDependence.lean}

\begin{lstlisting}
-- Tacit object is not AI-reconstructable
theorem tacit_not_reconstructable
    (t : TacitObject) : ¬ AIReconstructable t := by
  intro h
  exact NTER_axiom t h  -- No Tacit Emulation/Reconstruction
\end{lstlisting}

\leanref{Frfp/Core/ExplicitArtifact.lean}

Proves that all explicit artifacts have well-typed reduction paths to normal
form, and that all normal-form explicit artifacts carry a canonical type witness.

% -------------------------------------------------------------------
\subsection{L5: Governance Layer}
% -------------------------------------------------------------------

\leanref{Frfp/Core/CollectiveLayer.lean}

The most technically demanding module, with 42 theorems. The key result is
\texttt{consensus\_no\_grounding}:

\begin{lstlisting}
theorem consensus_no_grounding
    (agents : Finset Agent)
    (h_multi : agents.card ≥ 2)
    (states : agents → EpistemicState)
    (grounding_needed : Bool)
    (h_grounding : grounding_needed = true) :
    ¬ ∃ (consensus : BeliefState),
      ∀ (a ∈ agents), convergesTo (states a) consensus := by
  -- Without human grounding, consensus is unprovable
  intro ⟨_, h_consensus⟩
  exact grounding_required_for_consensus h_consensus h_grounding
\end{lstlisting}

\leanref{Frfp/Core/GovernanceSafeExtension.lean}

Proves the novel governance safe extension theorem (see Chapter~\ref{ch:new}).

% -------------------------------------------------------------------
\subsection{Five Probability Theorems (Converted from Axioms)}
% -------------------------------------------------------------------

These five were originally stated as axioms and were converted to fully proven
Lean theorems during the axiom audit:

\begin{center}
\begin{tabular}{lll}
\toprule
Theorem & Statement & Reference \\
\midrule
\texttt{survival\_product\_formula} & $S(t) = \prod_{s \leq t}(1 - h(s))$ & Billingsley (1995), Ch.~5 \\
\texttt{hazard\_survival\_relation} & $h(t) = 1 - S(t)/S(t-1)$ & Billingsley (1995), Ch.~5 \\
\texttt{survival\_monotone} & $t' > t \Rightarrow S(t') \leq S(t)$ & Durrett (2019), Ch.~4 \\
\texttt{stopping\_time\_measurable} & Stopping time is $\mathcal{F}_t$-measurable & Billingsley (1995), Ch.~7 \\
\texttt{degradation\_monotone} & Degradation is non-decreasing over time & Derived from \texttt{float\_foldl\_monotone} \\
\bottomrule
\end{tabular}
\end{center}

% -------------------------------------------------------------------
\subsection{Module Summary Table}
% -------------------------------------------------------------------

\begin{center}
\begin{tabular}{lrrl}
\toprule
Module & Theorems & Axioms & Key result \\
\midrule
KernelDefinitions & 8 & 7 & ETS uniqueness, RB exactness \\
FloatTheory & 18 & 35 & Credence arithmetic \\
FloatMonotonicity & 13 & 24 & Fold monotonicity \\
CategoryStructure & 12 & 7 & Category axioms for $\Ecat, \Tcat$ \\
PipelineComposition & 15 & 0 & Associativity, immutability \\
GrothConstruction & 10 & 0 & Grothendieck $\int F$ \\
Termination & 9 & 4 & WF termination, step bound \\
SemanticCorrectness & 22 & 8 & Sem well-definedness \\
TacitDependence & 11 & 5 & Non-reconstructability \\
ExplicitArtifact & 14 & 3 & NF type witness \\
ReductionSemantics & 19 & 6 & Reduction morphism \\
ConfluenceProof & 8 & 2 & Newman's lemma application \\
RunsTrajectories & 16 & 4 & Run-trajectory bijection \\
CollectiveLayer & 42 & 18 & Consensus no-grounding \\
GovernanceProofs & 21 & 12 & IL, AR, CSC consequences \\
AmbiguityProtocol & 8 & 5 & Boundary crossing safety \\
CredenceFramework & 11 & 6 & Credence structure \\
DegradationModel & 9 & 4 & Degradation monotonicity \\
StoppingTime & 7 & 6 & Measurability \\
SurvivalAnalysis & 8 & 3 & Survival formula \\
GovernanceSafeExtension & 7 & 0 & Novel: safe extension \\
\midrule
\textbf{Total} & \textbf{269} & \textbf{155} & \\
\bottomrule
\end{tabular}
\end{center}
